\paragraph{Changing what the network learns.}
The offer comparison prompted a paired follow-up on the same cohort. We replaced only the target in the 2025 training data, retaining each applicant's enrolled program when it appeared on the regular applications, or the outside class otherwise. This changes 45,450 of 189,982 training labels. We used the same row masks, features, shuffling, two architectures, seeds, optimizer, epochs and 2026 test features. The outside class describes missing on-list enrollment and includes registration through other routes. We selected no epoch or scoring rule using the test results.

All six H200 runs completed. With enrollment as the training target, the rank-blind and rank-informed menus covered 119,414 and 130,305 of the 136,751 linked registrations. The informed menus thus covered 10,891 more actual destinations. Replayed DA placed 158,498 and 158,357 applicants, respectively, a loss of 141. The informed model kept 612 more original full-list recipients placed (1,239 versus 627 losses) but brought in 753 fewer previously unmatched applicants (1,040 versus 287 entrants). Training on registrations therefore preserved the within-architecture sign pair on this cohort. It also reveals a distinction between objectives: versus the offer-trained informed model, the enrollment-trained informed model covered 637 more registered destinations, covered 79 fewer initial offers, and placed 176 fewer applicants. The individual-level records and predictions stay out of the shared archive.
